\section{Products and the continuous totient counting scale}
\label{sec:products}

The recurrence enters \eqref{eq:counting-scale} multiplicatively, so
an additive estimate should be transported through the product with
its sign and its first correction intact. Define
\begin{equation}\label{eq:product-definitions}
 \delta_n=\frac{\rho^n E_n}{\gamma},\qquad
 P_m=\prod_{j=1}^m\frac{g_j}{\gamma\rho^{-j}}
     =\prod_{j=1}^m(1-\delta_j).
\end{equation}
Positivity of $g_j$ and the exact moment representation give
$0<\delta_j<1$. Moreover,
\begin{equation}\label{eq:delta-geometric}
 0<\delta_{j+1}<\rho\delta_j,
\end{equation}
because $E_{j+1}<E_j$. Hence $P_m$ decreases to a strictly positive
constant
\begin{equation}\label{eq:product-constant}
 C_*:=\prod_{j=1}^\infty(1-\delta_j)\in(0,1).
\end{equation}
Strict positivity follows from convergence of $\sum\delta_j$ and the
fact that none of the individual factors vanishes.

\begin{theorem}[Sharp product tail]
\label{thm:product}
For every fixed $K\ge0$, set $n=m+1$ and $L=\log n$. Then
\begin{equation}\label{eq:product-asymptotic}
 \log\frac{P_m}{C_*}
 =\frac{\rho^n}{\gamma(1-\rho)n^2}
    \left\{\sum_{k=0}^K\frac{c_k}{L^{k+2}}
                     +O_K(L^{-K-3})\right\}.
\end{equation}
In particular,
\begin{equation}\label{eq:product-leading}
 P_m=C_*\left[1+
   \frac{\rho^{m+1}}{\gamma(1-\rho)(m+1)^2\log^2(m+1)}
       \left(1+\frac2{\log(m+1)}+O(\log^{-2}(m+1))\right)
        \right].
\end{equation}
\end{theorem}

\begin{proof}
Write
\[
 T_m=\sum_{j=m+1}^\infty\delta_j,\qquad
 R_m=\sum_{j=m+1}^\infty[-\log(1-\delta_j)-\delta_j].
\]
Then $\log(P_m/C_*)=T_m+R_m$. The elementary power series for
$-\log(1-x)$ and \eqref{eq:delta-geometric} give
\begin{equation}\label{eq:log-remainder-bound}
 0\le R_m\le
 \frac{\delta_n^2}{2(1-\delta_n)(1-\rho^2)},\qquad n=m+1.
\end{equation}
Tonelli's theorem applied to the positive moment representation gives
the exact identity
\begin{equation}\label{eq:product-tail-integral}
 T_m=\frac{\rho^n}{\gamma}
       \int_0^1\frac{t^{n-1}\ell(t)}{1-\rho t}\,dt.
\end{equation}
Put $\beta=\rho/(1-\rho)$ and
\begin{equation}\label{eq:D-def}
 D_{n,r}=(-1)^r\Delta^r E_n
        =\int_0^1t^{n-1}(1-t)^r\ell(t)\,dt.
\end{equation}
Since $(1-\rho t)^{-1}=(1-\rho)^{-1}
[1-\beta(1-t)/(1+\beta(1-t))]$,
\begin{equation}\label{eq:tail-first-sandwich}
 \frac{\rho^n}{\gamma(1-\rho)}(E_n-\beta D_{n,1})
 \le T_m\le\frac{\rho^n E_n}{\gamma(1-\rho)}.
\end{equation}
The same endpoint argument that gives the main theorem yields
$D_{n,1}=O(n^{-3}\log^{-2}n)$; a proof with all coefficients is
given in Section~\ref{sec:additional-consequences}. Thus replacing
$(1-\rho t)^{-1}$ by $(1-\rho)^{-1}$ costs an extra factor $1/n$.
For every fixed $K$, this error and \eqref{eq:log-remainder-bound}
are absorbed by the last term of \eqref{eq:product-asymptotic}.
Exponentiating proves \eqref{eq:product-leading}.
\end{proof}

\subsection{Finite, computable two-sided bounds for the product constant}

The integral proof gives more than an asymptotic notation. For $n=m+1$,
let
\begin{align}\label{eq:product-finite-bounds}
 U_m&=\frac{\delta_n}{1-\rho}
       +\frac{\delta_n^2}{2(1-\delta_n)(1-\rho^2)},\nonumber\\
 L_m&=\max\left\{\delta_n,
   \frac{\rho^n}{\gamma(1-\rho)}(E_n-\beta D_{n,1})\right\}.
\end{align}
Every $m\ge0$ satisfies
\begin{equation}\label{eq:C-enclosure}
 \boxed{\qquad P_m e^{-U_m}\le C_*\le P_m e^{-L_m}.\qquad}
\end{equation}
The additional lower bound $T_m\ge\delta_n$ used here follows directly
from its defining positive sum. These formulas involve just finitely
many recurrence terms and the two scalar constants $\rho,\gamma$.
They become certified numerical enclosures when those inputs are
enclosed by outward-rounded interval arithmetic. The supplied
high-precision computations are diagnostics, and are not described
as interval certificates.

There is a systematic refinement. The exact finite geometric identity
gives, for every $r\ge1$,
\begin{align}\label{eq:alternating-tail}
 T_m={}&\frac{\rho^n}{\gamma(1-\rho)}
       \sum_{j=0}^{r-1}(-\beta)^jD_{n,j}
       +\mathcal R_{m,r},\nonumber\\
 0<&(-1)^r\mathcal R_{m,r}
       \le\frac{\rho^n\beta^r}{\gamma(1-\rho)}D_{n,r}.
\end{align}
No convergence of an infinite geometric series with ratio $\beta$ is
being used: $\beta>1$ is allowed. The signed remainder is the exact
integral with factor $(1+\beta(1-t))^{-1}$. For each fixed $r$,
$D_{n,r}=O_r(n^{-r-2}\log^{-2}n)$, so successive bounds resolve
successive powers of $1/n$.

\subsection{Effect on the counting scale}

From the definition of $P_m$ one obtains the exact algebraic identity
\begin{equation}\label{eq:counting-scale-product}
 \mathcal G_m(B)
 =\frac{B^m\rho^{m(m+1)/2}}{m!\gamma^m C_*}
       \exp\left(-\log\frac{P_m}{C_*}\right).
\end{equation}
It follows that the approximation obtained by replacing $P_m$ with
$C_*$ is an upper approximation. Its relative correction has the
sharp size
\begin{equation}\label{eq:counting-correction}
 \frac{\rho^{m+1}}{
   \gamma(1-\rho)(m+1)^2\log^2(m+1)}.
\end{equation}
This improves the product-tail estimate $O(\rho^m)$ obtainable from
a bounded absolute recurrence error by the factor
$m^{-2}\log^{-2}m$, and determines its leading constant and sign.
It is a statement about the scale \eqref{eq:counting-scale}. It does
not identify or eliminate independent arithmetic errors in a totient
counting argument.

\section{A renewal process and inversion of its convergence rate}
\label{sec:renewal-inversion}

Define a probability distribution on the positive integers by
\begin{equation}\label{eq:renewal-law}
 p_j=a_j\rho^j\quad(j\ge1).
\end{equation}
Its total mass is $A(\rho)=1$ and its mean is
$\sum jp_j=1/\gamma$. For independent positive integer-valued
increments with this law, let $u_n$ be the probability that a renewal
occurs at time $n$, with $u_0=1$. Splitting according to the first
increment gives
$u_n=\sum_{j=1}^n p_ju_{n-j}$, so
\begin{equation}\label{eq:renewal-u}
 u_n=\rho^n g_n,\qquad
 \gamma-u_n=\rho^n E_n\quad(n\ge1).
\end{equation}
Thus the renewal probabilities increase strictly to $\gamma$ from
time one onward, and
\begin{equation}\label{eq:renewal-sharp}
 \gamma-u_n=
 \frac{\rho^n}{n^2\log^2n}
 \left(1+\frac2{\log n}-\frac{\pi^2}{2\log^2n}
                         +O(\log^{-3}n)\right).
\end{equation}
The exponential rate is exactly $\rho$:
$\lim(\gamma-u_n)^{1/n}=\rho$. No bound $O(r^n)$ with $r<\rho$
can hold. The extra polynomial and logarithmic factors describe the
behavior at this exact exponential rate.

The moment formula also extends the error to real indices:
\begin{equation}\label{eq:continuous-error}
 E(x)=\int_0^1t^{x-1}\ell(t)\,dt,
 \qquad \delta(x)=\frac{\rho^xE(x)}\gamma\qquad(x\ge1).
\end{equation}
It is continuous for $x\ge1$, smooth for $x>1$, and strictly decreasing
to zero. All the asymptotic proofs above remain valid for real
$x\to\infty$.

\begin{theorem}[Inverse error law]
\label{thm:inverse}
Let $0<\varepsilon<\delta(1)$, and let $x_\varepsilon$ be the
unique real solution of $\delta(x_\varepsilon)=\varepsilon$.
Write
\[
 a=\log(1/\rho),\qquad
 B=\log(1/\varepsilon),\qquad T=\log B,\qquad c=\log a.
\]
Then as $\varepsilon\downarrow0$,
\begin{align}\label{eq:inverse-expansion}
 ax_\varepsilon={}&B-2T-2\log T+2c-\log\gamma
             +\frac{2c+2}{T}\nonumber\\
 &+\frac{c^2+2c-\pi^2/2-2}{T^2}+O(T^{-3}).
\end{align}
The least positive integer index with relative renewal error at most
$\varepsilon$ is exactly $\lceil x_\varepsilon\rceil$.
\end{theorem}

\begin{proof}
The logarithm of \eqref{eq:main-log}, including its first three
coefficients, gives
\begin{equation}\label{eq:inverse-equation}
 B=ax+2\log x+2\log\log x+\log\gamma
       -\frac2{\log x}
       +\frac{\pi^2/2+2}{\log^2x}
       +O(\log^{-3}x).
\end{equation}
Its first-order balance implies $x\sim B/a$. Substituting the
right side of \eqref{eq:inverse-expansion}, divided by $a$, gives
\[
 \log x=T-c+O(T/B),\qquad
 \log\log x=\log T-\frac cT-\frac{c^2}{2T^2}
                 +O(T^{-3}+B^{-1}).
\]
Matching the coefficients of $1/T$ and $1/T^2$ in
\eqref{eq:inverse-equation} gives $2c+2$ and
$c^2+2c-\pi^2/2-2$, respectively. The discarded error $T/B$ is
$O(T^{-3})$. Finally the exact function
$-\log\delta(x)$ has derivative
\[
 a-\frac{E'(x)}{E(x)}\ge a>0,
\]
because $E'(x)=\int t^{x-1}\log t\,\ell(t)\,dt<0$.
Therefore a residual $O(T^{-3})$ in the logarithmic equation changes
its solution by at most $O(T^{-3})$. Strict monotonicity proves the
statement about the least integer index.
\end{proof}

The last sentence uses the exact real solution. Rounding a truncated
asymptotic approximation requires an enclosure if that approximation
is close to an integer. At the purely asymptotic level the integer
index has an unavoidable $O(1)$ rounding uncertainty.

\section{Finite differences and tails of the spectral remainder}
\label{sec:additional-consequences}

The moment structure permits several extensions without a new complex
contour argument. They also verify that the product estimates above
use a quantitatively controlled error.

\begin{proposition}[Every fixed finite-difference order]
\label{prop:difference-asymptotic}
For fixed integers $r,K\ge0$,
\begin{equation}\label{eq:difference-asymptotic}
 D_{n,r}=\frac1{n^{r+2}}\left\{
   \sum_{k=0}^K\frac{d_{r,k}}{(\log n)^{k+2}}
             +O_{r,K}((\log n)^{-K-3})\right\},
\end{equation}
where
\begin{equation}\label{eq:difference-coefficients}
 d_{r,k}=(k+1)![u^k]\,
       \frac{e^{\gE u}\prod_{j=1}^{r+1}(j+u)}{\Gamma(1-u)}.
\end{equation}
If $H_s=\sum_{j=1}^s j^{-1}$ and
$H_s^{(2)}=\sum_{j=1}^s j^{-2}$, the first terms are
\begin{align}\label{eq:difference-first}
 d_{r,0}&=(r+1)!,\nonumber\\
 d_{r,1}&=2(r+1)!H_{r+1},\nonumber\\
 d_{r,2}&=3(r+1)!\left(H_{r+1}^2-H_{r+1}^{(2)}-\zeta(2)\right).
\end{align}
\end{proposition}

\begin{proof}
In \eqref{eq:D-def}, put $t=e^{-s}$. The local integrand is the
one for $E_n$ multiplied by $(1-e^{-s})^r=s^r(1+O_r(s))$.
The leading logarithmic sector is consequently $n^{-r-2}
F_{r+2}(\log n)$, while the discarded powers of $s$ cost an extra
factor $1/n$. The all-orders inverse-logarithm proof applies with
$y^{r+1}e^{-y}$ in place of $ye^{-y}$. The reflection identity for
$\Gamma$ simplifies the coefficient generator to
\eqref{eq:difference-coefficients}. Expanding its logarithm through
degree two gives \eqref{eq:difference-first}.
\end{proof}

Two tails deserve separate attention. They follow from exact summation
under the positive integral:
\begin{align}\label{eq:tail-integrals}
 \sum_{j>N}E_j&=\int_0^1\frac{t^N\ell(t)}{1-t}\,dt,\nonumber\\
 \sum_{j>N}jE_j&=\int_0^1
  \left(\frac{(N+1)t^N}{1-t}+\frac{t^{N+1}}{(1-t)^2}\right)
                    \ell(t)\,dt.
\end{align}
In particular,
\begin{align}\label{eq:tails-leading}
 \sum_{j>N}E_j
 &=\frac1{N\log^2N}
    \left(1-\frac{\pi^2}{2\log^2N}+O(\log^{-3}N)\right),\nonumber\\
 \sum_{j>N}jE_j
 &=\frac1{\log N}+\frac1{\log^2N}
           -\frac{\pi^2}{6\log^3N}+O(\log^{-4}N).
\end{align}
For the first formula the coefficient of $1/\log N$ inside the
parentheses vanishes: the generator is
$e^{\gE u}/\Gamma(1-u)$, without the factor $1+u$.
Alternatively, integrate the expansion of $E_x$ and bound the
sum--integral difference by the first term, using monotonicity.
For the second formula, termwise integration of
$x^{-1}(\log x)^{-2}$ and its corrections gives the displayed
constants. To justify the sum--integral error in the weighted case,
put $f(x)=xE(x)$. The moment formula gives $E'(x)\le0$, and
integration by parts, using $xE(x)\to0$, yields
\[
 \int_N^\infty|f'(x)|\,dx
 \le NE(N)+2\int_N^\infty E(x)\,dx
 =O(N^{-1}\log^{-2}N).
\]
Comparing $f(j)$ with its integral over $[j-1,j]$ bounds the
sum--integral error by this total variation. This does not require
differentiating an asymptotic error term.
Thus the weighted sum rule \eqref{eq:main-sums} converges only
logarithmically, even though $g_n$ itself grows exponentially.
